\section{强化学习中的延迟}
\label{sec:why_delay_ns}

如前所述，\gls{rl}及其\gls{mdp}模型是序贯决策的优质框架。然而，如读者所知，模型是对真实过程的简化，有其自身局限。审视\gls{rl}实际应用的文献可揭示这些局限所在。\cite{mahmood2018setting} 发现，在真实机器人应用中，从一系列候选障碍中，学习的两大主要障碍是动作空间的选择和状态观测或动作执行中的延迟。这些延迟通过向智能体提供过时信息或延迟智能体对环境的控制来影响智能体-环境交互。这是支持延迟研究的有力论据。

我们认为延迟是任何序贯问题固有的。人类面临的最普遍延迟类型是大脑处理来自身体传感器信息所需的时间。例如，运动信息从光感受器传递到大脑高级视觉区域所需时间估计约为 100 ms \cite{de1991vernier}。在\gls{rl}中，此延迟类似于传感器采集时间导致的延迟。延迟也可能来自智能体根据环境观测计算下一个动作所需的时间。例如，使用更复杂的策略（如\emph{深度神经网络（deep neural network）}）会加剧延迟，这在主动流动控制中已有观察到\cite{ren2021applying,mao2022active}。

延迟自然出现的另一个例子是交易。在实践中，信息收集和交易传输的速度极其重要。信息通过光纤走不同路径，导致在不同地点的异步到达时间。这些延迟可被利用以获利\cite{lewis2014flash}。此外，通过抽象掉市场上的其他智能体，单个交易智能体的回报受到延迟的负面影响\cite{wilcox1993effect}。使用\gls{rl}智能体的外汇（FX）实验证实了引入延迟后的性能下降\cite{liotet2022delayed}。延迟影响的其他例子包括带半挂车的拖拉机紧急制动时间\cite{bayan2010brake}；并行计算中计算节点可用信息\cite{hannah2018unbounded}；机械臂控制\cite{mahmood2018setting}。

即使环境看似同步运行（如围棋），延迟仍隐藏在某处。在关于 AlphaGo \cite{silver2016mastering}（与世界冠军李世石对弈的算法）开发历程的纪录片中，AlphaGo 花费了超过 30 秒计算第一步。围棋比赛总时间有限，过长的动作计算可能有害。在此例中延迟并未造成损害，因为 AlphaGo 仍有充足时间并最终战胜了李世石。

在仿真中，延迟可以人为去除，但这可能带走重要的真实感。在 Atari 2600 等游戏中，延迟显著影响\gls{rl}智能体的性能\cite{firoiu2018human}。一个简单解决方案是在智能体选择动作时暂停环境以消除延迟。然而，\cite{firoiu2018human} 认为，由于大多数游戏以 60Hz 显示图像，且研究发现某些人群的平均反应时间约为 250 ms \cite{jain2015comparative}，不为人工智能体建模延迟使其相对于人类获得不公平优势。

许多相关领域也考虑延迟问题。例如，在\emph{控制理论（control theory）}中\cite{chung1995time, dugard1998stability, niculescu2001delay, gu2003survey, richard2003time, fridman2014introduction}，延迟对系统稳定性有复杂影响\cite{kolmanovskii1999delay}。该领域特别考虑了拥塞控制问题，其中数据包从多个源发送到网络\cite{jacobson1988congestion}。在此问题中，延迟自然产生，因为数据包从不同源走不同路径\cite{witsenhausen1971separation, altman1999congestion,ait2003stability}。在\emph{在线学习（online learning）}中，接收延迟反馈会减慢学习过程\cite{joulani2013online,pike2018bandits,lancewicki2021stochastic}。在\emph{实时（real-time）}\gls{rl}中，智能体和环境之间的交互不再是同步的——即环境"等待"智能体选择动作——而是变成异步的。实际上，智能体必须在环境持续演变的同时计算下一个动作，且无法保证计算出的动作对新的环境状态仍然相关。有趣的是，该子领域已被证明等价于延迟\gls{rl}\cite[Theorem~1]{ramstedt2019real}。因此，该领域的兴趣\cite{hester2013texplore,caarls2015parallel,travnik2018reactive,ramstedt2019real,xiao2020thinking}也论证了延迟研究的价值。

基于所有上述原因，我们认为延迟是\gls{rl}实际应用中的核心研究方向。此外，在传统框架中引入延迟所引发的理论问题也具有研究价值，因为它们与\gls{rl}的许多子领域相关，如\glsfirst{pomdp}或实时\gls{rl}。
